\section{Discussion}
\label{sec:discussion}

\paragraph{Twelve to one is too far for diagnosis.} Our central result is negative and, we think, useful: under a criterion fixed in advance, no superclass diagnosis can be screened from lead I, and reduced configurations with two or six leads replace the 12-lead ECG only for ST/T change with the chest leads. The average loss for lead I (0.09 AUROC) is large, but the conclusion matrix shows where it comes from. The pattern is consistent with the spatial coverage of each configuration. Lead I views the heart's electrical axis from the lateral direction only, so inferior changes (seen mainly in II, III and aVF) and anterior or septal changes (seen mainly in V1--V4) are largely invisible to it; MI, which is defined by such regional changes, lost most. Hypertrophy criteria rely heavily on precordial voltages, and HYP was better preserved by the chest leads (0.883) than by the limb leads (0.855). ST/T changes are often most visible in the precordial leads, which is in line with STTC surviving the chest-lead configuration almost unchanged. These anatomical explanations are hypotheses consistent with the per-class results, not something our experiments test directly.

\paragraph{The bottleneck is information, not data or training.} Three independent observations point the same way. Supervised transfer from 12 leads did not improve any reduced-lead model after correction; adding ten times more training data improved the lead-I model by only 0.023 AUROC, whereas a 12-lead model trained on 10\% of the data still beat it; and different architectures converged to similar lead-I performance (0.824--0.838). Better optimisation of single-lead models is therefore unlikely to close the gap; additional information---a second lead, a chest lead or repeated recordings---is needed.

\paragraph{What a wearable can do.} Triage is the more realistic use. With lead I, models still separated normal from abnormal ECGs with AUROC 0.88 on PTB-XL and 0.92 on the external cohort, but their sensitivity at 90\% specificity (0.71 internally) fell short of what we required for a stand-alone screen. Such a model could support, but not replace, the decision to obtain a 12-lead ECG. Two of our findings matter for deployment. First, the lead-agnostic LA-ResNet handled every lead configuration with one set of weights at a cost of 0.003--0.010 AUROC, which is attractive for devices whose available leads vary. Second, inter-model disagreement identified records with roughly twice the error rate, which could be used to refer uncertain recordings for human review.

\paragraph{Operating point and criterion.} The second condition of the criterion (sensitivity $\geq0.80$ at 90\% specificity) is strict: even 12-lead models reached only about 0.74 for MI and 0.71 for HYP. We kept the criterion unchanged after seeing the results; relaxing it would change the verdicts for some cells but not for lead I, where the non-inferiority condition already fails by a wide margin (\cref{fig:conclusion}). The sensitivity analysis on the margin (0.01 and 0.03) did not produce any additional \emph{screen} cell.

\paragraph{Self-supervision and Transformers.} ECG-ViT was the weakest architecture with full labels and gained most from masked-signal pre-training with 10\% labels, consistent with the view that Transformers, which have weak inductive biases, benefit from pre-training. The gain was modest, and our pre-training had not converged within 100 epochs, so a longer or larger pre-training stage may help more. We did not compare masked modelling with other self-supervised objectives such as contrastive learning~\cite{mehari2022ssl}.

\paragraph{External validation.} Frozen models kept their discrimination on records from another country and hospital system, which suggests that what they learned on PTB-XL is not specific to its recording conditions. The higher external AUROC must not be over-interpreted. The external target was restricted by design (only strictly sinus-rhythm negatives and mapped diagnoses), PTB-XL keeps diagnostic statements of any likelihood and thus contains label noise that lowers internal AUROC, and the external positives were dominated by ST/T changes. When borderline voltage-only findings were counted as abnormal, the external AUROC fell to the internal level. Thresholds also shifted, so any deployment would need local recalibration.

\paragraph{Comparison with previous work.} Our 12-lead results (0.915--0.925) lie within the PTB-XL superclass benchmark (0.921--0.930)~\cite{strodthoff2021benchmark}, and our feature baseline (0.863) is close to its wavelet-feature model (0.874). Davarmanesh et al.~\cite{davarmanesh2024leadI} reported an AUC of 0.92 for MI versus normal records from lead I. Our lead-I MI AUROC was 0.80 for the primary approach, but the tasks differ: we classify MI against all other records, including other abnormalities, in a multi-label setting, so the numbers should not be compared directly.

\paragraph{Limitations.}
\begin{itemize}
  \item \emph{Single training source and simulated wearables.} All models were trained on PTB-XL at 100\,Hz. ``Lead I'' is lead I of a clinical recording, not a smartwatch recording, which differs in electrode placement, contact quality, motion artefacts and duration. The noise experiment used synthetic white noise rather than real wearable artefacts.
  \item \emph{Label definitions.} Superclass labels include low-likelihood statements, and the external mapping from SNOMED codes to superclasses is approximate; external results varied with the mapping.
  \item \emph{Statistical scope.} Ablations and label-efficiency curves used one seed; subgroup analyses have no confidence intervals; ECG-ViT pre-training had not converged.
  \item \emph{Fairness of comparison.} All models shared one training recipe without per-model tuning. This makes comparisons fair but means that no model is at its individual best, and the architectures differ in parameter count (0.8--4.1\,M).
  \item \emph{Interpretability.} Grad-CAM and attention maps are coarse and illustrative; they are not validated clinical explanations.
\end{itemize}

\paragraph{Future work.} Natural next steps are to evaluate on genuine smartwatch recordings paired with 12-lead ECGs, to test which \emph{additional} single lead (for example a chest lead obtainable by placing the watch on the chest) recovers most information, to tune the operating point for a two-stage triage pathway, and to extend self-supervised pre-training to larger unlabelled ECG collections.
